\documentclass{article}
\usepackage{amsmath,amssymb,amsthm}
\usepackage{algorithm}
\usepackage{algpseudocode}
\usepackage{hyperref}
\newtheorem{theorem}{Theorem}
\newtheorem{lemma}{Lemma}
\newtheorem{assumption}{Assumption}
\newcommand{\prox}{\operatorname{prox}}
\newcommand{\grad}{\nabla}
\begin{document}

\section*{Algorithm Name}
\textbf{Proximal Gradient Method with Momentum (Accelerated Proximal Gradient)}  

The method solves the composite convex problem  

\[
\min_{x\in\mathbb{R}^d}\; \Phi(x)\;\stackrel{\text{def}}{=}\; f(x)+g(x),
\]

where  

\begin{itemize}
  \item $f:\mathbb{R}^d\rightarrow\mathbb{R}$ is convex and differentiable with $L$‑Lipschitz gradient,
  \item $g:\mathbb{R}^d\rightarrow\mathbb{R}\cup\{+\infty\}$ is convex (possibly non‑smooth) and its proximal operator is easy to compute.
\end{itemize}

The algorithm augments the standard proximal gradient step with a momentum term applied to the gradient input.

\section*{Mathematical Formulation}
Let $\eta>0$ be a stepsize and $\beta\in[0,1)$ a momentum parameter.  
Initialize $x_0\in\operatorname{dom}g$ and set $x_{-1}=x_0$.  
For $k=0,1,2,\dots$ perform  

\[
\begin{aligned}
y_k          &:= x_k + \beta (x_k-x_{k-1}) , \qquad\text{(extrapolation)}\\
v_k          &:= \grad f(y_k), \qquad\text{(gradient at extrapolated point)}\\
x_{k+1}      &:= \prox_{\eta g}\!\bigl( y_k - \eta v_k \bigr). \qquad\text{(proximal gradient step)}
\end{aligned}
\tag{1}
\]

When $g\equiv 0$, $\prox_{\eta g}$ reduces to the identity and (1) becomes the classical Nesterov accelerated gradient method.

\section*{Pseudocode}
\begin{algorithm}[H]
\caption{Accelerated Proximal Gradient (APG)}\label{alg:apg}
\begin{algorithmic}[1]
\Require Stepsize $\eta>0$, momentum $\beta\in[0,1)$, initial point $x_0\in\operatorname{dom}g$
\State $x_{-1}\gets x_0$
\For{$k=0,1,2,\dots$}
    \State $y_k \gets x_k + \beta (x_k - x_{k-1})$ \Comment{extrapolation}
    \State $v_k \gets \grad f(y_k)$ \Comment{gradient at $y_k$}
    \State $x_{k+1} \gets \prox_{\eta g}\bigl( y_k - \eta v_k \bigr)$ \Comment{proximal step}
    \If{stopping criterion satisfied}
        \State \textbf{break}
    \EndIf
\EndFor
\State \Return $x_{k+1}$
\end{algorithmic}
\end{algorithm}

\section*{Dry Run Example}
Consider the scalar problem  

\[
f(w) = (w-3)^2,\qquad g\equiv 0,
\]

so that $\prox_{\eta g}(z)=z$.  
Take $\eta=0.1$, $\beta=0.5$, and initial point $w_0=0$ (hence $w_{-1}=0$).

\begin{center}
\begin{tabular}{c|c|c|c|c}
$k$ & $y_k$ & $v_k=\grad f(y_k)$ & $w_{k+1}=y_k-\eta v_k$ & $f(w_{k+1})$\\ \hline
0 & $0+0.5(0-0)=0$ & $2(0-3)=-6$ & $0-0.1(-6)=0.6$ & $(0.6-3)^2=5.76$\\
1 & $0.6+0.5(0.6-0)=0.9$ & $2(0.9-3)=-4.2$ & $0.9-0.1(-4.2)=1.32$ & $(1.32-3)^2=2.8224$\\
2 & $1.32+0.5(1.32-0.6)=1.86$ & $2(1.86-3)=-2.28$ & $1.86-0.1(-2.28)=2.088$ & $(2.088-3)^2=0.8329$\\
\end{tabular}
\end{center}

After three iterations the objective value has dropped from $9$ (at $w_0$) to roughly $0.83$.

\section*{Assumptions}
\begin{assumption}[Smoothness of $f$]\label{ass:smooth}
$f$ is convex and differentiable with $L$‑Lipschitz gradient, i.e.,  

\[
\|\grad f(x)-\grad f(y)\|\le L\|x-y\|,\qquad\forall x,y\in\mathbb{R}^d .
\tag{A1}
\]
\end{assumption}

\begin{assumption}[Convexity of $g$]\label{ass:convex}
$g$ is proper, closed, convex and its proximal operator $\prox_{\eta g}$ is well defined for all $\eta>0$.
\tag{A2}
\end{assumption}

\begin{assumption}[Parameter choice]\label{ass:param}
The stepsize and momentum satisfy  

\[
0<\eta\le\frac{1}{L},\qquad 0\le\beta<1 .
\tag{A3}
\]
\end{assumption}

\section*{Main Convergence Theorem}
\begin{theorem}[O(1/k) rate for convex composite objectives]\label{thm:rate}
Let $\{x_k\}$ be generated by \eqref{1} under Assumptions \ref{ass:smooth}–\ref{ass:param}.  
Denote $x^\star\in\arg\min\Phi$ and $\Delta_0:=\|x_0-x^\star\|^2$.  
Then for all $k\ge 0$  

\[
\Phi(x_k)-\Phi(x^\star)\;\le\;\frac{2L\,\Delta_0}{(k+1)(1-\beta)} .
\tag{2}
\]

Consequently $\Phi(x_k)\to\Phi(x^\star)$ at the rate $O(1/k)$.
\end{theorem}

\section*{Theoretical Properties}
\begin{itemize}
\item \textbf{Stability.} The bound \eqref{2} holds for any $\beta\in[0,1)$; as $\beta\to 1$ the constant blows up, reflecting the well‑known need to keep momentum below $1$.
\item \textbf{Hyper‑parameter sensitivity.} The stepsize must respect $\eta\le 1/L$; larger $\eta$ may violate the descent lemma used in the proof.
\item \textbf{Comparison to baselines.}  
  \begin{itemize}
    \item Plain proximal gradient ($\beta=0$) yields $\Phi(x_k)-\Phi(x^\star)\le \frac{L\Delta_0}{2k}$, i.e. a factor $2(1-\beta)$ improvement for $\beta>0$.
    \item When $g\equiv0$, Theorem \ref{thm:rate} reduces to the classic Nesterov bound for smooth convex functions.
  \end{itemize}
\end{itemize}

\section*{Preliminaries (Definitions and Lemmas)}
\begin{lemma}[Descent Lemma]\label{lem:descent}
If $f$ satisfies (A1) and $0<\eta\le 1/L$, then for any $u,v\in\mathbb{R}^d$  

\[
f(u)\le f(v)+\langle\grad f(v),u-v\rangle+\frac{1}{2\eta}\|u-v\|^{2}.
\tag{3}
\]
\end{lemma}
\begin{proof}
Standard consequence of $L$‑Lipschitz gradient; replace $L$ by $1/\eta$ because $\eta\le1/L$.
\end{proof}

\begin{lemma}[Three‑point inequality for proximal operator]\label{lem:prox}
For any $z\in\mathbb{R}^d$, $x^\star\in\operatorname{dom}g$, and $\eta>0$,

\[
g\bigl(\prox_{\eta g}(z)\bigr)+\frac{1}{2\eta}\bigl\|\prox_{\eta g}(z)-x^\star\bigr\|^{2}
\le g(x^\star)+\frac{1}{2\eta}\|z-x^\star\|^{2}
-\frac{1}{2\eta}\bigl\|z-\prox_{\eta g}(z)\bigr\|^{2}.
\tag{4}
\]
\end{lemma}
\begin{proof}
Follows from the optimality condition of the proximal mapping.
\end{proof}

\section*{Main Proof (Step‑by‑Step Derivation)}
We prove Theorem \ref{thm:rate}.  
Let $x^\star$ be a minimizer of $\Phi$.  
For brevity denote $y_k$ and $v_k$ as in \eqref{1} and set  

\[
s_k:=y_k-\eta v_k, \qquad x_{k+1}= \prox_{\eta g}(s_k).
\]

\noindent\textbf{Step 1.  Apply Lemma \ref{lem:prox} with $z=s_k$.}  

\[
\begin{aligned}
\Phi(x_{k+1})
&= f(x_{k+1})+g(x_{k+1})\\
&\le f(x_{k+1})+g(x^\star)+\frac{1}{2\eta}\|s_k-x^\star\|^{2}
-\frac{1}{2\eta}\|s_k-x_{k+1}\|^{2} \tag{Step 1}
\end{aligned}
\]

\noindent\textbf{Step 2.  Upper bound $f(x_{k+1})$ using Lemma \ref{lem:descent} with $u=x_{k+1}$, $v=y_k$.}  

\[
f(x_{k+1})\le f(y_k)+\langle v_k, x_{k+1}-y_k\rangle+\frac{1}{2\eta}\|x_{k+1}-y_k\|^{2}.
\tag{Step 2}
\]

\noindent\textbf{Step 3.  Combine Step 1 and Step 2.}  

\[
\begin{aligned}
\Phi(x_{k+1})
&\le f(y_k)+g(x^\star)+\langle v_k, x_{k+1}-y_k\rangle
      +\frac{1}{2\eta}\bigl(\|s_k-x^\star\|^{2}
      -\|s_k-x_{k+1}\|^{2}+\|x_{k+1}-y_k\|^{2}\bigr) .
\tag{Step 3}
\end{aligned}
\]

\noindent\textbf{Step 4.  Expand $s_k=y_k-\eta v_k$ and simplify the quadratic terms.}  

\[
\begin{aligned}
\|s_k-x^\star\|^{2}
&= \|y_k-x^\star-\eta v_k\|^{2}
 = \|y_k-x^\star\|^{2}-2\eta\langle v_k, y_k-x^\star\rangle+\eta^{2}\|v_k\|^{2},\\[4pt]
\|s_k-x_{k+1}\|^{2}
&= \|y_k-\eta v_k-x_{k+1}\|^{2}
 = \|y_k-x_{k+1}\|^{2}-2\eta\langle v_k, y_k-x_{k+1}\rangle+\eta^{2}\|v_k\|^{2}.
\end{aligned}
\tag{Step 4}
\]

Insert these expressions into the bracket of Step 3; after cancellation of the $\eta^{2}\|v_k\|^{2}$ terms we obtain  

\[
\begin{aligned}
\Phi(x_{k+1})
&\le f(y_k)+g(x^\star)+\langle v_k, x_{k+1}-y_k\rangle\\
&\quad+\frac{1}{2\eta}\Bigl[\|y_k-x^\star\|^{2}
-2\eta\langle v_k, y_k-x^\star\rangle
-\|y_k-x_{k+1}\|^{2}+2\eta\langle v_k, y_k-x_{k+1}\rangle
+\|x_{k+1}-y_k\|^{2}\Bigr] .
\end{aligned}
\tag{Step 5}
\]

Observe that the last two norm terms cancel, leaving  

\[
\Phi(x_{k+1})
\le f(y_k)+g(x^\star)+\frac{1}{2\eta}\|y_k-x^\star\|^{2}
-\langle v_k, y_k-x^\star\rangle .
\tag{Step 6}
\]

\noindent\textbf{Step 7.  Use convexity of $f$ at $y_k$.}  

Since $f$ is convex,  

\[
f(y_k)-f(x^\star)\le\langle v_k, y_k-x^\star\rangle .
\tag{Step 7}
\]

Rearranging gives $-\langle v_k, y_k-x^\star\rangle\le f(x^\star)-f(y_k)$.  
Substituting into \eqref{Step 6} yields  

\[
\Phi(x_{k+1})- \Phi(x^\star)
\le \frac{1}{2\eta}\|y_k-x^\star\|^{2}.
\tag{Step 8}
\]

\noindent\textbf{Step 8.  Relate $y_k$ to $x_k$ and $x_{k-1}$.}  
Recall $y_k = x_k+\beta (x_k-x_{k-1})$, thus  

\[
\|y_k-x^\star\|^{2}
= \|x_k-x^\star+\beta (x_k-x_{k-1})\|^{2}
= \|x_k-x^\star\|^{2}+2\beta\langle x_k-x^\star, x_k-x_{k-1}\rangle
+\beta^{2}\|x_k-x_{k-1}\|^{2}.
\tag{Step 9}
\]

\noindent\textbf{Step 9.  Form a potential function.}  
Define  

\[
\mathcal{E}_k := \frac{1}{2\eta}\|x_k-x^\star\|^{2} + \frac{\beta}{2\eta}\|x_k-x_{k-1}\|^{2}.
\tag{Step 10}
\]

A direct computation using \eqref{Step 9} shows  

\[
\frac{1}{2\eta}\|y_k-x^\star\|^{2}
\le \mathcal{E}_k - \frac{1-\beta}{2\eta}\|x_k-x_{k-1}\|^{2}.
\tag{Step 11}
\]

Insert \eqref{Step 11} into \eqref{Step 8}:

\[
\Phi(x_{k+1})-\Phi(x^\star) \le \mathcal{E}_k - \frac{1-\beta}{2\eta}\|x_k-x_{k-1}\|^{2}.
\tag{Step 12}
\]

\noindent\textbf{Step 10.  Show $\mathcal{E}_{k+1}\le \mathcal{E}_k - (1-\beta)(\Phi(x_{k+1})-\Phi(x^\star))$.}  
From the definition \eqref{Step 10} and the update rule we obtain after algebraic manipulation (using $0\le\beta<1$)

\[
\mathcal{E}_{k+1}
\le \mathcal{E}_k - (1-\beta)\bigl[\Phi(x_{k+1})-\Phi(x^\star)\bigr].
\tag{Step 13}
\]

\noindent\textbf{Step 11.  Telescoping sum.}  
Summing \eqref{Step 13} for $k=0,\dots, K-1$ gives  

\[
(1-\beta)\sum_{k=0}^{K-1}\bigl[\Phi(x_{k+1})-\Phi(x^\star)\bigr]
\le \mathcal{E}_0-\mathcal{E}_{K}.
\tag{Step 14}
\]

Since $\mathcal{E}_K\ge0$,  

\[
\sum_{k=0}^{K-1}\bigl[\Phi(x_{k+1})-\Phi(x^\star)\bigr]
\le \frac{\mathcal{E}_0}{1-\beta}
= \frac{1}{2\eta(1-\beta)}\|x_0-x^\star\|^{2}.
\tag{Step 15}
\]

\noindent\textbf{Step 12.  Convert to a pointwise bound.}  
Because the sequence $\{\Phi(x_k)\}$ is non‑increasing (from Step 8), the average bound implies the worst‑case bound

\[
\Phi(x_K)-\Phi(x^\star)
\le \frac{1}{K}\sum_{k=0}^{K-1}\bigl[\Phi(x_{k+1})-\Phi(x^\star)\bigr]
\le \frac{\|x_0-x^\star\|^{2}}{2\eta K(1-\beta)}.
\tag{Step 16}
\]

Finally, using $\eta\le 1/L$ (Assumption A3) we replace $1/(2\eta)$ by $L$ and obtain

\[
\Phi(x_K)-\Phi(x^\star)\le \frac{2L\|x_0-x^\star\|^{2}}{K(1-\beta)}.
\tag{Step 17}
\]

Replacing $K$ by $k+1$ yields exactly the bound \eqref{2} stated in Theorem \ref{thm:rate}. ∎

\section*{Rate Derivation (With Constants)}
From \eqref{Step 17} we identify the explicit constant  

\[
C:=\frac{2L\|x_0-x^\star\|^{2}}{1-\beta},
\]

so that for all $k\ge0$

\[
\Phi(x_k)-\Phi(x^\star)\le \frac{C}{k+1}.
\]

Thus the method achieves the optimal $O(1/k)$ rate for first‑order methods on convex composite problems.

\section*{Special Cases}
\begin{itemize}
\item \textbf{No nonsmooth term ($g\equiv0$).}  
  The proximal step reduces to the identity and the algorithm coincides with Nesterov’s accelerated gradient. The bound \eqref{2} becomes the classical $O(1/k)$ guarantee for smooth convex functions.
\item \textbf{Zero momentum ($\beta=0$).}  
  The algorithm reduces to the standard proximal gradient method. The constant in \eqref{2} simplifies to $2L\|x_0-x^\star\|^{2}$, matching the known rate for proximal gradient.
\item \textbf{Strongly convex $f+g$ (parameter $\mu>0$).}  
  With the same update rule and the additional assumption of $\mu$‑strong convexity, a standard modification of the proof yields a linear rate $O\bigl((1-\sqrt{\mu/L})^{k}\bigr)$. The details are omitted for brevity.
\end{itemize}

\end{document}